\section{Introduction}

\emph{Whew -- I've finished the work, now I only have to write the report.} Sound familar? Many students put off writing until too late and don't leave enough time to make a good job of the report. This is a serious mistake because the report typically accounts for a large fraction of the assessment. Start early -- most people find report-writing difficult. Think about the report throughout your project, keep track of references and collect material to illustrate the report.

\subsection{Mechanical aspects}

The length of the body of the report will be specified in the instructions for the project. Extra material may be provided in appendices but this material should be for reference only: you cannot assume that the reader will study it. In other words, do not put vital points in an appendix.

You will probably think that the report is too short but this is deliberate. Most reports are submitted to busy managers, who do not have time to read lengthy documents. It is important to learn how to pick out the vital points and write a concise report with maximum impact.

Reports should be word-processed and firmly bound. Use A4 paper and a clear typeface such as 12-point Times, number the pages and leave margins of at least 25\,mm all round. Follow this layout (this document breaks some of the rules to keep it compact).
\begin{itemize}
\item
The front cover should show the title of the project with your name(s) and matriculation number(s) and the name of your supervisor(s).
\item
Put the abstract on the next page. It should be about 100--250 words and gives a brief summary of the report including the background and aims of the project, the principal results and conclusions.
\item
The next page should show the table of contents. 
\item
The body of the report should be divided into numbered sections, each starting on a new page. Figures (diagrams, plots or photographs) and tables need captions and should be numbered.
\end{itemize}

\subsection{What goes into the introduction?}

Explain the background to the project and the reasons why your particular piece of work was considered worthwhile. This leads to the \textbf{aims} of the project: what you are trying to achieve. Be specific.